\section{Introduction}\label{sec:intro}

Wearable proximity sensors record face-to-face contacts between people at a resolution of seconds \citep{cattuto2010,stehle2011}, and the resulting temporal networks \citep{holme2012} have become a standard empirical basis for models of social contact and epidemic spreading.
Their most robust statistical feature is the broad distribution of contact durations and of the intervals between successive contacts of the same pair, far from the exponential distributions of a memoryless process.
The prevailing explanation is an intrinsic memory of the individuals: the probability to leave a contact or a group decreases with the time already spent in it, as in the reinforcement models of \citet{zhao2011} and \citet{starnini2013}, or ties are reinforced by repeated interaction, as in activity-driven networks with memory \citep{karsai2014}.
Such aging makes the process non-Markovian in the network state.

An alternative explanation is that the process is Markovian, but in a richer state than the set of current contacts.
A contact embedded in a group (with common contacts), between members of the same class, between friends, or of a student who currently has many other contacts, may have very different rates of ending than an isolated contact between strangers, and a mixture of such rates, changing as groups form and dissolve, also produces broad durations.
Traces of recent interactions that decay over time add memory to such a description while keeping it Markovian in an extended state, since the traces are part of the state.
The two explanations have different consequences: intrinsic aging puts the memory into the individuals, network context puts it into the configuration around them, and only the latter is captured by models of network structure.

Here we quantify how much of the memory is network context.
We specify a hierarchy of Markovian models with increasing context and memory, and a competing non-Markovian aging model, for the contacts of \Students\ students of a high school \citep{mastrandrea2015}, fit all models by maximum likelihood on three days, and compare them on two held-out days by their log-likelihood and by time-rescaling residuals \citep{brown2002}.
The Markovian models are rule sets for \tricl, a simulator and inference engine for typed multi-relational network dynamics, which computes the exact log-likelihood and its gradient for each rule set, including rates that decay between events.

We find that memory dominates and that network context explains only a small part of it, but that a Markovian model with decaying traces on several time scales outperforms power-law aging.
We also find that all fitted models with network feedback are dynamically unstable, a form of model degeneracy that a pair closure of the generated mean-field equations predicts and a standard closure misses.
Section~\ref{sec:data} describes the data, Sec.~\ref{sec:models} the models, Sec.~\ref{sec:inference} inference and evaluation, Sec.~\ref{sec:results} the results, Sec.~\ref{sec:closure} the stability analysis and closures, and Sec.~\ref{sec:discussion} the conclusions.
